\documentclass[10pt,a4paper]{amsart}
\usepackage[margin=19mm]{geometry}
\usepackage{amsmath,amssymb,mathtools}
\usepackage[scr=rsfs,cal=euler]{mathalfa}
\usepackage{xfrac}
\usepackage[unicode,hidelinks,bookmarks=false]{hyperref}
\newcommand{\ie}{i.e.\ }
\newcommand{\cf}{cf.\ }
\newcommand{\resp}{respectively\ }
\newcommand{\Subsection}{\subsection}
\providecommand{\nobreakdash}{\nobreak\hbox{-}}
\setlength{\emergencystretch}{2em}
\multlinegap0pt
\usepackage{amssymb}
\newtheorem{lemma}{Lemma}
\newcommand{\K}{\mathcal{K}}
\title{When Arcs Extend Uniquely: A Higher-Dimensional Generalization of Barlotti’s Result}
\author{Tim L. Alderson}
\date{Source: arXiv:2511.06193v1 (CC BY 4.0)}
\begin{document}
\maketitle

\noindent\emph{Source and licence.} When Arcs Extend Uniquely: A Higher-Dimensional Generalization of Barlotti’s Result. Version arXiv:2511.06193v1 (CC BY 4.0), distributed by arXiv under the Creative Commons Attribution 4.0 International licence, http://creativecommons.org/licenses/by/4.0/, as recorded in the arXiv licence metadata for this exact version and shown on the abstract page https://arxiv.org/abs/2511.06193v1. Full licence notice: \texttt{licenses/arxiv-2511.06193-CC-BY-4.0.txt}.\par\smallskip

\noindent\textit{Editorial scope.} Counted unit: the article's lemma result (source0.tex lines 147--149). The article's complete original proof of it is delivered with it (source0.tex lines 151--174, one \texttt{proof} environment of 24 lines). No mathematics is written, repaired or re-derived: the statement and the proof are the article's own lines, in the article's own notation, and no retained line is substituted or re-flowed.  Retained uncounted as the source-local context the proof needs: lines 124--126; lines 133--135.  Nothing was omitted from any retained span, so the package carries no omission record.  The retained lines cite 1 works, whose entries are transcribed verbatim from the article's own bibliography source in the e-print and delivered with short editorial labels recorded in the item's reference mapping.  No retained line contains a figure environment, an image-inclusion command or drawing code, so no image is delivered and the package declares no asset dependency.  Editorial formatting only, in addition: an AMS \texttt{amsart} preamble that restates the article's own declarations and macros used by the retained lines, namely the environments \texttt{lemma} on separate counters, together with the standard packages those lines need.  The retained environments print in this document as Lemma 1 in order of appearance, whereas the article prints the same items with the numbering of its own full document.  The complete native source of the article as distributed in the arXiv e-print for this exact version is bundled unchanged as \texttt{sources/188716/source0.tex}.
\begin{lemma}[\cite{MR0083141}]\label{lem Barlotti Extension k = 3}
	If \( \K \) is an \(((s+1)(q+1), s+2)\)-arc in PG\((2, q)\), with \( 0 < s < q - 2 \), and \( (s+2) \mid q \), then there exists a unique point \( P \) not incident with any secant of \( \K \), and consequently incident with all tangents to \( \K \). Hence, \( \K \) admits a unique extension to a maximal \(((s+1)(q+1)+1, s+2)\)-arc.
\end{lemma}

\begin{lemma}\label{lem: long implies projective and cap}
	Let \( \K \) be an \((n, k+s-1)\)-(multi-)arc in PG\((k-1, q)\), with \( k \ge 3 \). If \( n > s(q+1) + k - 1 \), then each $(k-3)$-flat intersects $\K$ in at most $k-2$ points. Consequently,  \( \K \) must be a set (i.e., not a multi-arc).
\end{lemma}

\begin{lemma}\label{lem: main result}
	Let \( \K \) be an \((n, k + s - 1)\)-arc in PG\((k - 1, q)\), with \( k \ge 3 \), \( 0 < s < q - 2 \), and \( (s+2) \mid q \). If \( n = (s+1)(q+1) + k - 3 \), then there exists a unique point \( X \in \text{PG}(k - 1, q) \) that is incident with all tangents to \( \K \), and not incident with any secant of \( \K \). Consequently, \( \K \) is not complete and admits a unique extension to a maximal \((n + 1, k + s - 1)\)-arc.
\end{lemma}

\begin{proof}
	By Lemma~\ref{lem: long implies projective and cap}, the arc \( \K \) must be a set (i.e., not a multi-arc). We proceed by induction on the projective dimension \( k - 1 \).
	
	\smallskip
	\noindent \text{Base case \( k = 3 \):} This is precisely Lemma~\ref{lem Barlotti Extension k = 3}, which establishes the existence of a unique point not lying on any secant of \( \K \) and incident with all tangents.
	
	\smallskip
	%\noindent \textbf{Inductive step:} 
	Assume the result holds for PG\((k-2, q)\). Let \( \K = \{P_1, P_2, \ldots, P_n\} \subset \Pi = \text{PG}(k-1, q) \), with \( n = (s+1)(q+1) + k - 3 \), and \( k \ge 4 \). For each point \( P_i \in \K \), consider the quotient geometry \( \Pi^*_i = \Pi / P_i \cong \text{PG}(k - 2, q) \), and let \( \K_i^* \) be the image of \( \K \setminus \{P_i\} \) under this quotient.
	
	By construction, \( \K_i^* \) is an arc of size \( (s+1)(q+1) + k - 4 \) in PG\((k - 2, q)\), satisfying the hypotheses of the inductive step. Thus, there exists a unique line \( \ell_i \subset \Pi \) through \( P_i \), lying entirely in the tangents to \( \K \) at \( P_i \), and meeting every secant through \( P_i \) only at \( P_i \) itself.
	
	\smallskip
	\noindent \textit{Claim:} The lines \( \ell_i \) are pairwise distinct. Suppose \( \ell_i = \ell_j \) for some \( i \ne j \). In the quotient at \( P_i \), the image \( P_j^* \in \K_i^* \) corresponding to $\ell_j$ is incident with no secants. Thus $\K_i^*$ admits an extension to a multi-arc, violating Lemma~\ref{lem: long implies projective and cap}. 
	
	\smallskip
	\noindent \text{ \( k = 4 \):} Counting shows that for $i\ne j$, the points $P_i,P_j$ are incident with an unique tangent of $\K$, denoted $\pi_{ij}$, so $\ell_i$ and $\ell_j$ meet (in $\pi_{ij}$). Take any $P_a \notin \pi_{12}$. The line $\ell_a$ similarly meets both $\ell_1$ and $\ell_2$, and must meet $\pi_{12}$ in an unique point, say $X$.  Thus, each $\ell_i$ in $\pi_{12}$ meets $\ell_a$ in $X$, and each $\ell_i\notin \pi_{1,2}$ meets $\ell_1$ and $\ell_2$ in $X$, so all \( \ell_i \) are concurrent at \( X \).
	
	\smallskip
	\noindent \text{ \( k \ge 5 \):} With appeal to Lemma \ref{lem: long implies projective and cap}, for any pair \( P_i, P_j \in \K \), the quotient at \( \langle P_i, P_j \rangle \) yields a maximal arc in PG\((k-3, q)\). The inductive hypothesis implies that  there exists an unique plane $\pi_{ij}$ through $\langle P_i,P_j\rangle $ which is contained in each tangent of $\K$ through $\langle P_i,P_j\rangle $, and which meets each secant of $\K$ through $\langle P_i,P_j\rangle $ precisely in $\langle P_i,P_j\rangle $.  Note that by definition, $\ell_i,\ell_j \subseteq \pi_{ij}$ for all $i,j$.	
	
 
	We now claim that the set \( \{\ell_i\} \) is copunctal.  Indeed, consider $\ell_1, \ell_2$, and $\ell_3$.  Since $\ell_1,\ell_2 \subseteq \pi_{12}$, they meet at a point, $X$. If $P_3\notin \pi_{12}$ then $\ell_1=\pi_{13}\cap\pi_{12}$, so $\ell_3\cap \pi_{12}=Y\in \ell_1$.  Likewise, by considering $\pi_{23}$ we obtain $\ell_3\cap \pi_{12}=Y\in \ell_2$, giving $Y=X$.  Thus $X\in \ell_1,\ell_2,\ell_3$. Hence, all \( \ell_i \) pass through the common point \( X \), which is incident with all tangents and avoids all secants of \( \K \), completing the proof.
\end{proof}

\begin{thebibliography}{99}
\bibitem[MR0083]{MR0083141} Adriano Barlotti. \newblock Sui {$\{k;n\}$}-archi di un piano lineare finito. \newblock {\em Boll. Un. Mat. Ital. (3)}, 11:553--556, 1956.
\end{thebibliography}
\end{document}
